\chapter{Fundamentos teóricos}
\label{chap:fundamentos}

La seguridad en el desarrollo de software depende no solo de la lógica de programación, sino también de las decisiones adoptadas durante la compilación. En este capítulo se describen las principales medidas de seguridad ofrecidas por distintos compiladores y cómo contribuyen a mitigar vulnerabilidades comunes. Se presentan primero las técnicas generales, disponibles en varios compiladores, y luego se detallan aquellas específicas de \emph{Clang} y \emph{Rustc}. Estas medidas constituyen la base teórica sobre la que se estructura el análisis de este trabajo.

\section{Medidas de seguridad evaluadas}

En este trabajo se analizan distintas opciones de compilación que
permiten reforzar la seguridad de los programas [26]. A continuación, se
describen de forma general las principales medidas evaluadas.

\begin{enumerate}
\def\labelenumi{\arabic{enumi}.}
\item
  \textbf{Ejecución independiente de la posición (del inglés, \emph{Position Independent Execution} o PIE)}.
  Permite que el programa se cargue en una dirección de memoria distinta
  cada vez que se ejecuta, lo que dificulta ciertos ataques basados en
  direcciones fijas de memoria. Esta característica se puede activar o
  desactivar mediante opciones específicas del compilador [1].
\item
  \textbf{Protección contra errores en la pila (\emph{stack protector})}.
  Añade medidas de seguridad para evitar que se modifique
  accidentalmente o de forma malintencionada la información que una
  función guarda temporalmente en la pila. Se puede ajustar el nivel de
  protección según las necesidades o incluso desactivarla si no se considera
  oprtuno [2].
\item
  \textbf{Fortificación de funciones estándar (\emph{Fortify source})}.
  Reemplaza funciones comunes de la biblioteca estándar (como copiar
  cadenas de texto) por versiones más seguras que realizan
  comprobaciones adicionales para evitar errores como por ejemplo desbordamientos.
  Se puede ajustar el nivel de fortificación o desactivarla por
  completo.
\item
  \textbf{Pila no ejecutable (\emph{Execstack})}.
  Evita que la parte de memoria destinada a la pila pueda ejecutarse. 
  Esta medida impide que un atacante inserte y ejecute código
  malicioso colocado en esa zona. También puede activarse o desactivarse
  durante la compilación.
\item
  \textbf{Protección de la tabla de enlaces (del inglés, \emph{Relocation Read-Only} o RELRO)}.
  Refuerza ciertas estructuras internas que se utilizan durante la carga del
  programa en memoria [27]. Estas estructuras son a veces objetivo de
  ataques y esta opción evita que puedan ser modificadas una vez
  cargadas [3].
\item
  \textbf{Sanitizadores}.
  Son herramientas que permiten detectar distintos tipos de errores
  durante la ejecución del programa. Entre los utilizados en este
  trabajo están:

  \begin{itemize}
  \item
    \textbf{AddressSanitizer}: identifica accesos incorrectos a memoria [4].
  \item
    \textbf{UndefinedBehaviorSanitizer}: detecta usos incorrectos del
    lenguaje que pueden provocar fallos difíciles de localizar [5].
  \item
    \textbf{MemorySanitizer}: avisa cuando se usa memoria sin haberla
    inicializado [6].
  \item
    \textbf{LeakSanitizer}: detecta si el programa pierde memoria sin
    liberarla (utilizado solo en Rust) [7].
  \end{itemize}
\item
  \textbf{Eliminación de redireccionamientos en llamadas a funciones externas (del inglés, \emph{Procedure Linkage Table} o PLT)}.
  Impide el uso de ciertas estructuras que facilitan la ejecución de
  funciones compartidas. Aunque puede mejorar el rendimiento,
  desactivar esta opción incrementa la seguridad ante ataques que modifican estos
  enlaces.
\end{enumerate}

Además de las opciones comunes descritas anteriormente, algunos
compiladores incorporan medidas de seguridad particulares que también se
han evaluado en este trabajo. Estas medidas se exponen a continuación.

\section{Medidas exclusivas de \emph{Clang}}

\begin{itemize}
\item
  \textbf{Pila segura (\emph{Safe Stack})}.
  Separa las variables locales sensibles del resto de la memoria de
  trabajo de una función, almacenándolas en una zona aislada y
  protegida. Esto reduce las posibilidades de que un atacante las
  manipule aprovechando errores como por ejemplo desbordamientos [8].
\item
  \textbf{Optimización en tiempo de enlace (del inglés, \emph{Link-Time Optimization} o
  LTO)}.
  Aunque no es una medida de seguridad directa, permite que el
  compilador analice y optimice el código completo al enlazar. Esto
  puede influir en cómo se aplican otras protecciones y en el
  comportamiento final del ejecutable.
\end{itemize}

\section{Medidas exclusivas de \emph{Rustc}}

\begin{itemize}
\item
  \textbf{Comprobación de desbordamientos aritméticos (\emph{Overflow checks})}.
  Verifica que las operaciones con enteros no sobrepasen los límites del
  tipo de dato. Si ocurre un desbordamiento, el programa lo detecta y lo
  informa [9].
\item
  \textbf{Aserciones de depuración (\emph{Debug assertions})}.
  Introduce comprobaciones adicionales durante la ejecución para ayudar
  a detectar errores lógicos o condiciones inesperadas en tiempo de
  desarrollo.
\item
  \textbf{Gestión de errores en tiempo de ejecución (\emph{panic})}.
  Permite configurar el comportamiento del programa ante fallos graves:
  puede detenerse inmediatamente (modo \emph{abort}) o intentar limpiar
  antes de salir (modo \emph{unwind}) [10].
\item
  \textbf{Optimización en tiempo de enlace (del inglés, \emph{Link Time Optimization} o LTO y \emph{Thin Link Time Optimization} o Thin LTO)}.
  Similar al caso de \emph{Clang}, \emph{Rustc} puede aplicar optimizaciones
  globales durante el proceso de enlace. La variante \emph{Thin LTO} es
  más rápida, pero con menor impacto.
\end{itemize}

\section{Relación con vulnerabilidades comunes}

Las medidas de seguridad presentadas en este capítulo están pensadas para reducir los riesgos frente a errores y ataques comunes en programas modernos. A continuación, se explican algunas amenazas frecuentes y cómo estas protecciones ayudan a evitarlas.

Los desbordamientos de búfer [11] ocurren cuando un programa escribe más datos de los que una variable puede almacenar, sobrescribiendo otras zonas de memoria. Esto puede usarse para tomar el control del programa. Técnicas como el \emph{stack protector}, el uso de pila no ejecutable y el mecanismo \emph{Safe Stack} ayudan a evitar este tipo de problemas al vigilar y proteger la memoria que usan las funciones.

Los ataques conocidos [25] como \emph{Return-Oriented Programming} (ROP) [12] [28] intentan ejecutar código dañino reutilizando partes del programa original. Para hacer esto, el atacante necesita saber exactamente dónde están ciertos fragmentos de código en memoria. Medidas como PIE, la protección de las estructuras internas del programa y la desactivación de ciertos mecanismos de salto (como PLT) hacen que el binario se cargue de forma diferente cada vez, lo que dificulta este tipo de ataque.

Los errores de memoria, como usar memoria ya liberada o no inicializada, o dejar sin liberar la memoria usada, pueden hacer que el programa falle o se comporte de forma inesperada. Los sanitizadores, como \emph{AddressSanitizer} o \emph{LeakSanitizer}, ayudan a encontrar estos problemas durante las pruebas del programa, haciendo que sean más fáciles de detectar y corregir.

Los errores en operaciones con números, como cuando un número entero se sale de su límite y ``vuelve a empezar desde cero", pueden causar fallos sutiles y difíciles de rastrear. En Rust, este tipo de errores se pueden detectar automáticamente si se activa la comprobación de desbordamientos, especialmente durante el desarrollo.

Algunos errores en el código, aunque no siempre provocan fallos visibles, pueden llevar a situaciones incorrectas o peligrosas si no se controlan. Herramientas como \emph{UndefinedBehaviorSanitizer} y las aserciones de depuración ayudan a identificar este tipo de comportamientos que el lenguaje considera ``indefinidos", como dividir por cero en algunas arquitecturas o acceder fuera de un arreglo.

Finalmente, los fallos graves que el programa no sabe manejar pueden provocar que se quede colgado o que actúe de forma errática. En Rust, el sistema de manejo de errores en tiempo de ejecución (\emph{panic}) permite decidir si se quiere detener el programa inmediatamente o intentar cerrarlo de forma ordenada.